\section{Conclusion}
\label{sec:10}
\label{sec:conclusion}

VoteChain-ZK extends the previously developed biometric-blockchain voting architecture
with a candidate-blind authorization and encrypted-ballot pathway. Biometric verification
authorizes release of a protected voting credential without receiving the voter's
candidate selection. A client-side privacy worker uses that credential to construct a
fixed-slot encrypted ballot and an election-bound Groth16 proof of eligibility, nullifier
consistency, and ballot validity. The voting contract verifies the proof, rejects reused
election-scoped nullifiers, and accumulates ciphertexts without learning the selected
candidate. Following voting, three of five trustees support aggregate decryption, with
their evidence checked externally before the contract records final totals and a
transcript commitment.

The controlled evaluation recorded 2,000 accepted encrypted-ballot transactions across 25
elections, with mean gas consumption of 1,314,307.123 per accepted ballot and mean local
submission-to-receipt latency of 760.640 ms. Candidate-count trials showed limited gas
variation within the fixed 16-slot circuit. Under the fixed depth-17 Merkle configuration,
eligible-set preparation costs increased with population size, whereas mean Groth16
proving time remained between 5.292 and 5.423 seconds. These observations characterize the
tested components and execution environment; they do not measure complete
biometric-to-receipt voting time or public-network throughput.

The resulting privacy and verifiability guarantees remain conditional on credential
custody, client integrity, trustworthy cryptographic setup, and trustee-share protection.
In particular, the custodian may link voter identities to election nullifiers, and the
contract does not verify trustee decryption proofs: independent checking requires the
complete external finalization transcript. Future work should evaluate end-to-end voting
sessions and concurrent public-network operation, establish independently managed trustee
custody and production-grade setup, and investigate stronger identity unlinkability and
expanded circuit configurations. The present results support a bounded implementation of
encrypted, proof-checked voting and externally auditable tally recovery, not a
production-election security claim.
